
\PassOptionsToPackage{hidelinks}{hyperref}
\documentclass[
  aps,
  prd,
  reprint,
  superscriptaddress,
  nofootinbib,
  floatfix,
  longbibliography
]{revtex4-2}

\usepackage[T1]{fontenc}
\usepackage[utf8]{inputenc}
\usepackage{amsmath,amssymb,amsthm,mathtools,bm}
\usepackage{booktabs,array,tabularx}
\usepackage{microtype}
\IfFileExists{orcidlink.sty}{\usepackage{orcidlink}}{\newcommand{\orcidlink}[1]{}}

\newtheorem{definition}{Definition}
\newtheorem{proposition}{Proposition}
\newtheorem{theorem}{Theorem}
\newtheorem{corollary}{Corollary}
\newtheorem{lemma}{Lemma}
\newtheorem{criterion}{Criterion}
\newtheorem{remark}{Remark}

\newcommand{\Sthree}{S^3}
\newcommand{\D}{\mathcal D}
\newcommand{\vis}{\mathrm{vis}}
\newcommand{\crit}{\mathrm{crit}}
\newcommand{\eff}{\mathrm{eff}}
\newcommand{\TT}{\mathrm{TT}}
\newcommand{\BY}{\mathrm{BY}}
\newcommand{\MS}{\mathrm{MS}}
\newcommand{\bigO}{\mathcal O}
\newcommand{\dd}{\mathrm d}
\newcommand{\Inv}{\mathcal I}
\newcommand{\N}{\mathbb N}
\newcommand{\R}{\mathbb R}
\newcommand{\Z}{\mathbb Z}

\begin{document}

\title{Supplemental material: Analytic invariant benchmarks for compact \texorpdfstring{\(S^3\)}{S3} residual cosmology}

\author{Andrei~Preece\,\orcidlink{0009-0004-5978-795X}}
\email{andre.preece@web.de}
\affiliation{Independent Researcher, Berlin, Germany}

\author{Boris~Batenin\,\orcidlink{0009-0000-7027-8583}}
\email{batenin@web.de}
\affiliation{Independent Researcher, Berlin, Germany}

\date{\today}

\begin{abstract}
This supplement gives analytic benchmark identities for the compact \(S^3\)
residual-invariant framework developed in the companion benchmark paper.  The word
``benchmark'' is essential: the results
below are exact reference cases and consistency checks, not a claim that all
solutions of the Einstein-matter system on \(I\times S^3\) are known.  We derive:
(i) the lapse-independent autonomous background system and its elementary fixed
points; (ii) the scalar heat trace, its correct ultraviolet and infrared
asymptotics, and \(\zeta_\Delta(0)\); (iii) admissibility conditions for isotropic
scalar covariance kernels; (iv) curvature-mimicking and constant-\(w\) residual
backgrounds; (v) acceleration, bounce, and distance criteria; and (vi) cautious
topological and group-theoretic constraints.  Several tempting but incorrect
statements are explicitly excluded: Brown--York and Misner--Sharp energies do not
coincide to all orders for a homogeneous density, the Bousso bound is not a
volume-scaling formula, the odd-dimensional Dirac index does not exclude all
gauge-coupled zero modes, and the four-dimensional Gauss--Bonnet identity does not
reduce to a constraint on \(\int R^2\) alone in a general FLRW history.
\end{abstract}

\maketitle

%=======================================================================
\section{Purpose and conventions}
\label{sec:purpose}
%=======================================================================

The companion benchmark paper starts from
\begin{equation}
  (M,g_{\mu\nu},\Sigma,U^\mu,T^\vis_{\mu\nu}),\qquad
  M\simeq I\times S^3,
\end{equation}
and defines the residual stress tensor
\begin{equation}
  T^\D_{\mu\nu}:=\frac{c^4}{8\pi G}G_{\mu\nu}-T^\vis_{\mu\nu}.
  \label{eq:residual}
\end{equation}
The present supplement records analytic identities that can serve as validation
tests for numerical or Bayesian implementations of that framework.
The external inputs are standard: FLRW dynamical systems~\cite{WainwrightEllis1997}, spectral theory on compact Riemannian manifolds~\cite{Berger1971,Rosenberg1997}, positive-definite kernels on spheres~\cite{Schoenberg1942}, quasi-local energies~\cite{BrownYork1993,MisnerSharp1964}, the covariant entropy bound~\cite{Bousso2002}, the Einstein constraint equations~\cite{Wald1984}, the Lichnerowicz bound~\cite{Lichnerowicz1958}, and Ray--Singer torsion~\cite{RaySinger1971}.

We use the closed FLRW metric
\begin{equation}
  ds^2=-N^2(t)c^2dt^2+a^2(t)R_0^2d\Omega_3^2,
\end{equation}
and distinguish coordinate-time and proper-time derivatives:
\begin{equation}
  f':=\frac{1}{N}\frac{df}{dt},\qquad H=\frac{a'}{a}.
  \label{eq:proper}
\end{equation}
The present-day curvature parameter is
\begin{equation}
  \Omega_K=-\frac{c^2}{H_0^2A_0^2}<0,\qquad A_0=a_0R_0,\qquad a_0=1 .
  \label{eq:OmegaK}
\end{equation}
Thus a curvature term in the Friedmann equation appears as \(\Omega_K a^{-2}\),
with \(\Omega_K<0\) for \(S^3\).

\begin{remark}
A closed topology is a global statement.  The late-time normalised curvature
density may be arbitrarily small while the spatial slice remains \(S^3\).
\end{remark}

%=======================================================================
\section{Autonomous background system}
\label{sec:phase}
%=======================================================================

Let the components scale as \(a^{-\alpha_i}\), with
\[
  \alpha_m=3,\qquad \alpha_r=4,\qquad
  \alpha_\D=3(1+w_\D),\qquad \alpha_K=2 .
\]
Define the normalised variables
\begin{equation}
  x_i:=\frac{\Omega_i a^{-\alpha_i}}{E^2(a)},\qquad
  E(a)=\frac{H(a)}{H_0}.
\end{equation}
For the curvature component, \(x_K<0\).  The Friedmann constraint is
\begin{equation}
  x_m+x_r+x_\D+x_K=1.
  \label{eq:xconstraint}
\end{equation}
The lapse-independent autonomous system is
\begin{equation}
  \frac{dx_i}{d\ln a}=x_i(Q-\alpha_i),
  \qquad
  Q:=\sum_j\alpha_j x_j .
  \label{eq:auto}
\end{equation}
The deceleration parameter is
\begin{equation}
  q=-1+\frac{Q}{2}
   =\frac12 x_m+x_r+\frac12(1+3w_\D)x_\D .
  \label{eq:q}
\end{equation}
The curvature variable drops out of \(q\).  This is not an approximation; curvature
contributes to \(H^2\) but not directly to \(a''/a\).

\begin{theorem}[Elementary fixed points and degeneracies]
\label{thm:fixed}
Assume constant \(w_\D\).  The isolated vertex fixed points of
Eq.~\eqref{eq:auto} are the single-component vertices whose scaling exponent
equals \(Q\):
\begin{align}
  R &: (x_r,x_m,x_\D,x_K)=(1,0,0,0),\quad q=1,\\
  M &: (0,1,0,0),\quad q=\frac12,\\
  D &: (0,0,1,0),\quad q=\frac12(1+3w_\D),\\
  K &: (0,0,0,1),\quad q=0 .
\end{align}
The \(K\) vertex is a formal fixed point of the projective system, but it is not a
physical point of a positively expanding \(S^3\) model because \(x_K<0\) for
\(S^3\).

If two components have equal scaling exponent, the corresponding edge is fixed.
In particular:
\begin{enumerate}
\item \(w_\D=0\): matter and residual dust form a fixed line.
\item \(w_\D=1/3\): radiation and residual radiation form a fixed line.
\item \(w_\D=-1/3\): residual stress and curvature form a fixed line with
      \(x_\D+x_K=1\).  For \(S^3\) this physical line has
      \(x_K<0\) and hence \(x_\D>1\), not \(0\le x_\D\le 1\).
\end{enumerate}
\end{theorem}

\begin{proof}
At a fixed point each component satisfies either \(x_i=0\) or \(Q=\alpha_i\).
If only one component is nonzero, the constraint gives the vertex.  If two or
more nonzero components have distinct \(\alpha_i\), the equations require
simultaneously \(Q=\alpha_i=\alpha_j\), impossible.  Thus nontrivial fixed
edges occur exactly when the scaling exponents coincide.
\end{proof}

\begin{proposition}[Linearised eigenvalues at the physical vertices]
\label{prop:eigs}
For a vertex with exponent \(\alpha_*\), perturbing by a component \(j\) gives
linear eigenvalue
\begin{equation}
  \lambda_j=\alpha_*-\alpha_j
  \label{eq:eig_general}
\end{equation}
with respect to \(d/d\ln a\).  Thus at radiation:
\[
  \lambda_m=1,\qquad \lambda_\D=1-3w_\D,\qquad \lambda_K=2.
\]
At matter:
\[
  \lambda_r=-1,\qquad \lambda_\D=-3w_\D,\qquad \lambda_K=1.
\]
At a cosmological-constant vertex \(w_\D=-1\):
\[
  \lambda_m=-3,\qquad \lambda_r=-4,\qquad \lambda_K=-2.
\]
Therefore the \(\Lambda\) vertex is the future attractor inside the expanding
physical sector whenever \(\Omega_\Lambda>0\) and no phantom or recollapse branch
intervenes.
\end{proposition}

\begin{proposition}[Einstein static benchmark]
\label{prop:ES}
For dust plus a cosmological constant in a closed FLRW model,
\[
  E^2=\Omega_m a^{-3}+\Omega_\Lambda-|\Omega_K|a^{-2}.
\]
The Einstein static solution satisfies \(H=0\) and \(a''=0\), hence
\begin{equation}
  a_{\rm ES}=\sqrt{\frac{|\Omega_K|}{3\Omega_\Lambda}},
  \qquad
  \Omega_m=2\Omega_\Lambda a_{\rm ES}^3 .
  \label{eq:ES}
\end{equation}
The physical radius is
\begin{equation}
  A_{\rm ES}=A_0a_{\rm ES}=\frac{c}{H_0\sqrt{3\Omega_\Lambda}}.
  \label{eq:AES}
\end{equation}
With the convention \(\Omega_\Lambda=\Lambda c^2/(3H_0^2)\), this becomes
\begin{equation}
  A_{\rm ES}=\frac{1}{\sqrt{\Lambda}}.
\end{equation}
The linearised equation is
\begin{equation}
  \delta a''=3H_0^2\Omega_\Lambda\,\delta a,
\end{equation}
so the static solution is unstable.
\end{proposition}

%=======================================================================
\section{Scalar spectral benchmarks}
\label{sec:spectral}
%=======================================================================

The scalar Laplacian on \(S_A^3\) has
\begin{equation}
  \lambda_n=\frac{n(n+2)}{A^2},\qquad d_n=(n+1)^2,\qquad n=0,1,\ldots .
  \label{eq:scalar_spec}
\end{equation}

\begin{proposition}[Exact scalar heat trace]
\label{prop:heat}
The scalar heat trace is
\begin{align}
  K(s)&=\mathrm{Tr}\,e^{-s(-\Delta)}
  =\sum_{n=0}^{\infty}(n+1)^2e^{-sn(n+2)/A^2}
  \nonumber\\
  &=e^{s/A^2}\sum_{m=1}^{\infty}m^2e^{-sm^2/A^2}.
  \label{eq:heat}
\end{align}
Poisson summation gives the exact dual representation
\begin{align}
  K(s)&=e^{s/A^2}\frac{A^3\sqrt{\pi}}{4s^{3/2}}
  \Bigg[
  1+2\sum_{k=1}^{\infty}
  \left(1-\frac{2\pi^2A^2k^2}{s}\right)
  \nonumber\\
  &\hspace{2.2cm}\times e^{-\pi^2A^2k^2/s}
  \Bigg].
  \label{eq:heat_poisson}
\end{align}
\end{proposition}

\begin{corollary}[Correct ultraviolet and infrared limits]
\label{cor:UVIR}
As \(s\to0^+\),
\begin{align}
  K(s)&\sim
  \frac{V}{(4\pi s)^{3/2}}
  \left[1+\frac{\mathcal R^{(3)}}{6}s+\bigO(s^2/A^4)\right],
  \nonumber\\
  V&=2\pi^2A^3,\qquad \mathcal R^{(3)}=\frac{6}{A^2}.
  \label{eq:weyl}
\end{align}
As \(s\to\infty\),
\begin{equation}
  K(s)=1+4e^{-3s/A^2}+9e^{-8s/A^2}+\bigO(e^{-15s/A^2}).
  \label{eq:IR}
\end{equation}
The limiting constant \(1\) is the multiplicity-one scalar zero mode.  Equation~\eqref{eq:heat_poisson} is exact, but it is not an infrared asymptotic
series when truncated at fixed \(k\); large cancellations recover
Eq.~\eqref{eq:IR}.  Thus it is better described as a Poisson dual
representation, not as a literal T-duality.
\end{corollary}

\begin{theorem}[Scalar zeta value at zero]
\label{thm:zeta0}
For the nonzero scalar spectrum,
\begin{equation}
  \zeta_\Delta(z)=
  \sum_{n=1}^{\infty}(n+1)^2\left[\frac{n(n+2)}{A^2}\right]^{-z}
  =A^{2z}\sum_{m=2}^{\infty}\frac{m^2}{(m^2-1)^z},
  \label{eq:zeta}
\end{equation}
analytic continuation gives
\begin{equation}
  \zeta_\Delta(0)=-1 .
  \label{eq:zeta0}
\end{equation}
This is a scale-independent zeta invariant of the scalar Laplacian with its zero
mode removed.
\end{theorem}

\begin{proof}
For \(m\ge2\),
\[
  \frac{m^2}{(m^2-1)^z}=m^{2-2z}(1-m^{-2})^{-z}.
\]
Expanding in powers of \(m^{-2}\) and analytically continuing,
\[
  \sum_{m=2}^\infty \frac{m^2}{(m^2-1)^z}
  =
  \sum_{k=0}^\infty
  \binom{-z}{k}(-1)^k
  \left[\zeta_R(2z+2k-2)-1\right].
\]
At \(z=0\), only \(k=0\) survives and
\(\zeta_R(-2)-1=-1\).  Since \(A^0=1\), Eq.~\eqref{eq:zeta0} follows.
\end{proof}

\begin{corollary}[Determinant scaling]
\label{cor:det}
If \(\det'(-\Delta):=\exp[-\zeta_\Delta'(0)]\), then
\begin{equation}
  \log\det'(-\Delta_A)=2\log A+C_{S^3},
  \label{eq:det}
\end{equation}
where \(C_{S^3}\) is the unit-radius scalar constant.  Only the \(2\log A\)
scaling is used as a benchmark; the constant depends on the chosen operator and
normalisation.
\end{corollary}

\begin{proposition}[Scalar representation and Lichnerowicz equality]
\label{prop:SO4}
The isometry group of \(S_A^3\) is
\(\mathrm{SO}(4)\simeq(\mathrm{SU}(2)\times\mathrm{SU}(2))/\mathbb Z_2\).
Scalar harmonics of degree \(n\) transform in the \((n/2,n/2)\) representation,
whose dimension is \((n+1)^2\).  The first nonzero scalar eigenvalue is
\begin{equation}
  \lambda_1=\frac{3}{A^2}.
\end{equation}
Since \(\mathrm{Ric}=2A^{-2}g\), the Lichnerowicz bound in dimension three gives
\(\lambda_1\ge3/A^2\), and \(S^3\) saturates the bound.
\end{proposition}

\begin{remark}
Precise vector and transverse-traceless tensor spectra require fixing the
operator: rough Laplacian, Hodge Laplacian, Lichnerowicz operator, or a gauge-fixed
cosmological perturbation operator.  They should not be inferred by shifting the
scalar spectrum.  In the benchmark pipeline, these sectors must be recorded in an
operator ledger before numerical predictions are compared.
\end{remark}

%=======================================================================
\section{Admissible scalar covariance kernels}
\label{sec:cov}
%=======================================================================

Let
\begin{equation}
  C(x,y)=\sum_{n=0}^{\infty}A_n\,\Xi_{n,A}(x,y)
\end{equation}
be an isotropic scalar kernel, where the addition kernel satisfies
\begin{equation}
  \Xi_{n,A}(x,x)=\frac{(n+1)^2}{2\pi^2A^3}.
\end{equation}

\begin{theorem}[Admissibility of isotropic scalar covariance]
\label{thm:cov}
The kernel \(C\) is a positive semidefinite covariance kernel for a scalar
Gaussian field if
\begin{equation}
  A_n\ge0
  \label{eq:Anpos}
\end{equation}
and the point variance is finite:
\begin{equation}
  \sum_{n=0}^{\infty}A_n\,\Xi_{n,A}(x,x)
  =
  \frac{1}{2\pi^2A^3}
  \sum_{n=0}^{\infty}A_n(n+1)^2
  <\infty .
  \label{eq:variance}
\end{equation}
Conversely, every square-integrable isotropic scalar covariance has a spectral
expansion of this form with nonnegative coefficients.  This is the \(S^3\)
harmonic form of Schoenberg positivity.
\end{theorem}

\begin{corollary}[Benchmark spectra]
\label{cor:spectra}
The power law \(A_n=A_0(n+1)^{-\alpha}\) has finite point variance exactly when
\(\alpha>3\).  The exponential family
\begin{equation}
  A_n=A_0e^{-\beta n},\qquad \beta>0,
\end{equation}
is admissible, with
\begin{equation}
  \sigma^2=
  \frac{A_0}{2\pi^2A^3}
  \frac{e^\beta(e^\beta+1)}{(e^\beta-1)^3}.
\end{equation}
The scale-invariant analogue \(A_n\propto(n+1)^{-2}\) has divergent point variance
without an ultraviolet cutoff or additional smoothing.
\end{corollary}

%=======================================================================
\section{Residual background benchmarks}
\label{sec:bg}
%=======================================================================

\subsection{Curvature-mimicking residual stress}

\begin{theorem}[Exact \(w_\D=-1/3\) class]
\label{thm:wm13}
For \(w_\D=-1/3\),
\[
  \Omega_\D(a)=\Omega_{\D0}a^{-2}.
\]
In a model containing only this residual component and the \(S^3\) curvature term,
\begin{equation}
  E^2(a)=\frac{\Omega_{\D0}-|\Omega_K|}{a^2}.
  \label{eq:wm13}
\end{equation}
There are three cases:
\begin{enumerate}
\item If \(\Omega_{\D0}>|\Omega_K|\), then
\begin{equation}
  a(\tau)=a_i+H_0\sqrt{\Omega_{\D0}-|\Omega_K|}\,(\tau-\tau_i),
  \label{eq:coast}
\end{equation}
so \(q=0\).
\item If \(\Omega_{\D0}=|\Omega_K|\), then \(H=0\); the background is static
within this two-term model.
\item If \(\Omega_{\D0}<|\Omega_K|\), then no real expanding solution exists in
this two-term model.
\end{enumerate}
At background level the \(w_\D=-1/3\) residual is exactly degenerate with a shift
of the curvature coefficient.  Perturbations are needed to break the degeneracy.
\end{theorem}

\subsection{Constant-\texorpdfstring{\(w\)}{w} turning points}

\begin{proposition}[Turning point for a single residual plus curvature]
\label{prop:turning}
For constant \(w_\D\) and no matter or radiation,
\begin{equation}
  E^2(a)=\Omega_{\D0}a^{-3(1+w_\D)}-|\Omega_K|a^{-2}.
\end{equation}
If \(w_\D\ne -1/3\), the turning scale is
\begin{equation}
  a_t=
  \left(\frac{\Omega_{\D0}}{|\Omega_K|}\right)^{1/(1+3w_\D)} .
  \label{eq:at}
\end{equation}
For \(w_\D>-1/3\), \(a_t\) is a maximum expansion scale.  For
\(w_\D<-1/3\), \(a_t\) is a minimum scale on the allowed expanding branch.
The special value \(w_\D=-1/3\) is the curvature-mimicking case of
Theorem~\ref{thm:wm13}.
\end{proposition}

\begin{corollary}[Examples]
For \(w_\D=-1\),
\[
  a_t=\sqrt{\frac{|\Omega_K|}{\Omega_\Lambda}},
\]
the minimum radius of the pure closed de Sitter branch.  For \(w_\D=0\),
\[
  a_t=\frac{\Omega_{\D0}}{|\Omega_K|},
\]
the maximum radius of the closed dust model.  For \(w_\D=1/3\),
\[
  a_t=\sqrt{\frac{\Omega_{\D0}}{|\Omega_K|}} .
\]
\end{corollary}

\subsection{Acceleration and bounce criteria}

\begin{proposition}[Acceleration boundary]
\label{prop:acc}
The acceleration boundary is
\begin{equation}
  q=0
  \quad\Longleftrightarrow\quad
  \frac12x_m+x_r+\frac12(1+3w_\D)x_\D=0 .
  \label{eq:acc}
\end{equation}
For dust plus \(\Lambda\) and negligible radiation,
\begin{equation}
  a_{\rm acc}=\left(\frac{\Omega_m}{2\Omega_\Lambda}\right)^{1/3},
  \label{eq:aacc}
\end{equation}
independent of \(\Omega_K\).  The curvature parameter changes the distance-redshift
relation and the Friedmann constraint, but not the local sign of \(a''\) once the
component densities are specified.
\end{proposition}

\begin{proposition}[Regular FLRW bounce condition]
\label{prop:bounce}
A regular bounce at \(a=a_b>0\) requires
\begin{equation}
  H(a_b)=0,\qquad a''(a_b)>0 .
\end{equation}
Using the acceleration equation,
\begin{equation}
  \frac{a''}{a}=-\frac{4\pi G}{3c^2}
  \left(\epsilon_{\rm tot}+3p_{\rm tot}\right),
\end{equation}
the second condition requires
\begin{equation}
  \epsilon_{\rm tot}+3p_{\rm tot}<0
\end{equation}
at the bounce.  Thus violation of the perfect-fluid strong-energy combination is
necessary for a regular accelerating bounce in classical FLRW GR.  It is not,
by itself, sufficient: the Friedmann constraint must also vanish at \(a_b\).
\end{proposition}

\begin{remark}
The Hawking-Penrose singularity theorems have several hypotheses.  Violating
\(\epsilon+3p\ge0\) removes a key FLRW focusing condition, but it is not a complete
classification of singularity formation in general spacetimes.
\end{remark}

\subsection{BAO distance bound}

\begin{proposition}[Transverse-distance upper bound]
\label{prop:bao}
For a closed \(S^3\) FLRW geometry,
\begin{equation}
  D_M(z)=A_0\sin\!\left(\frac{\chi(z)}{A_0}\right),
  \qquad
  \chi(z)=\int_0^z\frac{c\,dz'}{H(z')}.
\end{equation}
Therefore
\begin{equation}
  D_M(z)\le A_0=\frac{c}{H_0\sqrt{|\Omega_K|}},
  \qquad
  \frac{D_M(z)}{r_d}\le
  \frac{c}{H_0r_d\sqrt{|\Omega_K|}} .
  \label{eq:bao_bound}
\end{equation}
The bound is geometric and model-independent once \(H_0,r_d,\Omega_K\) are fixed.
Whether it is observationally constraining depends on the redshift range and the
posterior for \(\Omega_K\).
\end{proposition}

%=======================================================================
\section{Interaction-current benchmark}
\label{sec:interaction}
%=======================================================================

\begin{proposition}[Scaling solution for a simple interaction current]
\label{prop:int}
Suppose
\begin{align}
  \epsilon'_\vis+3H(1+w_\vis)\epsilon_\vis &= Q_{\rm int},\\
  \epsilon'_\D+3H(1+w_\D)\epsilon_\D &= -Q_{\rm int},
\end{align}
with
\begin{equation}
  Q_{\rm int}=3H\Gamma(\epsilon_\D-\epsilon_\vis).
\end{equation}
For \(r:=\epsilon_\vis/\epsilon_\D\),
\begin{equation}
  \frac{1}{3H}\frac{r'}{r}
  =
  w_\D-w_\vis+\Gamma\left(\frac{1}{r}-r\right).
\end{equation}
A positive scaling fixed ratio satisfies
\begin{equation}
  \Gamma r_*^2+(w_\vis-w_\D)r_*-\Gamma=0,
\end{equation}
hence
\begin{equation}
  r_*=
  \frac{-(w_\vis-w_\D)+
  \sqrt{(w_\vis-w_\D)^2+4\Gamma^2}}{2\Gamma}
\end{equation}
for \(\Gamma>0\).  The earlier expression \(r_*=\Gamma^{-1}-1\) is not a
general solution; it can only arise after extra assumptions not present here.
\end{proposition}

%=======================================================================
\section{Quasi-local benchmarks}
\label{sec:quasi}
%=======================================================================

On a time-symmetric round \(S^3_A\) slice, let \(R_A=A\sin\chi\) be the areal
radius of a geodesic two-sphere.

\begin{proposition}[Brown--York energy on a round \(S^3\) slice]
\label{prop:BY}
With flat-space reference subtraction, the Brown--York energy is
\begin{equation}
  E_\BY(R_A)=
  \frac{c^4R_A}{G}
  \left[1-\sqrt{1-\frac{R_A^2}{A^2}}\right],
  \qquad 0\le R_A\le A .
  \label{eq:BY}
\end{equation}
The small-sphere expansion is
\begin{equation}
  E_\BY(R_A)=
  \frac{c^4}{2GA^2}R_A^3+
  \frac{c^4}{8GA^4}R_A^5+\bigO(R_A^7/A^6),
  \label{eq:BYsmall}
\end{equation}
and the equatorial value is
\begin{equation}
  E_\BY(A)=\frac{c^4A}{G}.
  \label{eq:BYmax}
\end{equation}
The associated mass is \(M_\BY=E_\BY/c^2\); one must not replace it by
\(c^4A/G^2\).
\end{proposition}

\begin{proposition}[Brown--York--Misner--Sharp comparison]
\label{prop:BYMS}
For homogeneous total energy density,
\begin{equation}
  E_\MS(R_A)=\frac{4\pi}{3}R_A^3\epsilon_{\rm tot}.
\end{equation}
Equality at a chosen radius \(R_A\) requires
\begin{equation}
  \epsilon_{\rm tot}(R_A)=
  \frac{3c^4}{4\pi G R_A^2}
  \left[1-\sqrt{1-\frac{R_A^2}{A^2}}\right].
  \label{eq:eps_BYMS}
\end{equation}
For a homogeneous density this cannot hold for all \(R_A\), because the right-hand
side depends on \(R_A\).  The leading small-sphere equality is obtained when
\begin{equation}
  \epsilon_{\rm tot}=\frac{3c^4}{8\pi GA^2}
  =\frac{c^4\mathcal R^{(3)}}{16\pi G},
\end{equation}
but the equality then fails at order \(R_A^5/A^4\).  Thus Brown--York and
Misner--Sharp energies are distinct benchmarks, not interchangeable definitions.
\end{proposition}

\begin{proposition}[Entropy bound: safe form]
\label{prop:entropy}
For a two-sphere of areal radius \(R_A\), the covariant entropy bound is bounded
by one quarter of its area in Planck units:
\begin{equation}
  S\le \frac{k_Bc^3}{4G\hbar}\,4\pi R_A^2
  =\frac{\pi k_Bc^3}{G\hbar}R_A^2 .
  \label{eq:entropy}
\end{equation}
The existence and orientation of a light-sheet are dynamical questions.  No
general \(S^3\) FLRW formula with an extra
\([1-\sqrt{1-R_A^2/A^2}]/(1+z)\) factor follows from the covariant entropy bound.
\end{proposition}

%=======================================================================
\section{Topological and index-theoretic benchmarks}
\label{sec:top}
%=======================================================================

\begin{proposition}[Untwisted Dirac benchmark]
\label{prop:dirac}
For the untwisted Dirac operator, using the Lichnerowicz formula~\cite{Lichnerowicz1958}, on a round \(S_A^3\),
\begin{equation}
  D^2=\nabla^*\nabla+\frac{\mathcal R^{(3)}}{4}
  =\nabla^*\nabla+\frac{3}{2A^2}.
\end{equation}
Hence the untwisted operator has no zero modes.  In odd dimension there is no
ordinary chiral Atiyah--Singer index of the even-dimensional type.  With a
nontrivial gauge connection, the twisted Lichnerowicz formula contains curvature
terms from the gauge field, and zero modes are not excluded by topology alone.
\end{proposition}

\begin{proposition}[Four-dimensional Gauss--Bonnet: correct use]
\label{prop:GB}
For a compact Riemannian four-manifold with boundary terms included,
\begin{equation}
  \chi(M)=\frac{1}{8\pi^2}\int_M
  \left(
  |W|^2-\frac12|\mathrm{Ric}_0|^2+\frac{R^2}{24}
  \right)dV
  +B_{\partial M}.
  \label{eq:GB}
\end{equation}
An FLRW spacetime is conformally flat, so \(W=0\), but the traceless Ricci term
generally does not vanish.  Therefore Gauss--Bonnet does not reduce, in a general
FLRW history, to a topological constraint on \(\int R^2dV\) alone.  For
\(I\times S^3\) with \(I\) a finite interval, \(\chi(I\times S^3)=0\); since the
manifold has boundary, the boundary term is essential.
\end{proposition}

\begin{proposition}[Closed geodesic periods]
\label{prop:geo}
On the round simply connected \(S_A^3\), every geodesic is closed with primitive
length
\begin{equation}
  L_1=2\pi A,\qquad L_k=2\pi kA .
\end{equation}
The exponential factors \(e^{-\pi^2A^2k^2/s}=e^{-L_k^2/(4s)}\) appearing in the
Poisson form of the heat trace encode these periods.  This is a useful spectral
geometric dictionary.  It should not be called a Selberg trace formula in the
hyperbolic sense, since the closed geodesics on \(S^3\) occur in continuous
families rather than as isolated hyperbolic orbits.
\end{proposition}

\begin{proposition}[Ray--Singer torsion benchmark]
\label{prop:torsion}
For the round \(S^3\) with the trivial real representation, the Ray--Singer
analytic torsion is \(T(S^3)=1\).  This is a topological benchmark for the
de Rham complex.  It does not impose a direct constraint on the residual stress
tensor unless an additional quantum effective-action construction is specified.
\end{proposition}

%=======================================================================
\section{Extended invariant ledger}
\label{sec:ledger}
%=======================================================================

The entries below extend the companion ledger.  Each entry is marked as exact,
conditional, or cautionary.

\begin{enumerate}
\setcounter{enumi}{84}

\item \(\Inv_{85}\): \(Q=\sum_i\alpha_i x_i\), exact autonomous source scalar.

\item \(\Inv_{86}\): \(q=-1+Q/2\), lapse-independent deceleration scalar.

\item \(\Inv_{87}\): vertex eigenvalues \(\lambda_j=\alpha_*-\alpha_j\), exact
linearised benchmark.

\item \(\Inv_{88}\): physical curvature condition \(x_K<0\) for \(S^3\), which
excludes the formal \(K\)-vertex from the expanding physical simplex.

\item \(\Inv_{89}\): \(a_{\rm ES}=\sqrt{|\Omega_K|/(3\Omega_\Lambda)}\), exact
Einstein-static scale factor for dust plus \(\Lambda\).

\item \(\Inv_{90}\): \(A_{\rm ES}=1/\sqrt{\Lambda}\) when
\(\Omega_\Lambda=\Lambda c^2/(3H_0^2)\), exact physical radius.

\item \(\Inv_{91}\): \(K(s)=e^{s/A^2}\sum_{m=1}^\infty m^2e^{-sm^2/A^2}\), exact
scalar heat trace.

\item \(\Inv_{92}\): Poisson representation Eq.~\eqref{eq:heat_poisson}, exact
dual heat-trace form.

\item \(\Inv_{93}\): ultraviolet Weyl coefficient
\(V(4\pi s)^{-3/2}\), exact leading local spectral density.

\item \(\Inv_{94}\): curvature heat coefficient
\(\mathcal R^{(3)}s/6=s/A^2\), exact first local correction.

\item \(\Inv_{95}\): infrared scalar heat limit \(K(s)\to1\), exact dominance
of the single constant scalar zero mode.

\item \(\Inv_{96}\): \(\zeta_\Delta(0)=-1\), exact scale-independent zeta value
for the nonzero scalar spectrum.

\item \(\Inv_{97}\): \(\log\det'(-\Delta_A)=2\log A+C_{S^3}\), exact determinant
scaling.

\item \(\Inv_{98}\): scalar Lichnerowicz equality \(\lambda_1=3/A^2\), exact.

\item \(\Inv_{99}\): covariance positivity \(A_n\ge0\), exact admissibility
condition.

\item \(\Inv_{100}\): finite point variance
\(\sum_nA_n(n+1)^2/(2\pi^2A^3)<\infty\), exact scalar covariance criterion.

\item \(\Inv_{101}\): exponential covariance benchmark
\(A_n=A_0e^{-\beta n}\), exact admissible family for \(\beta>0\).

\item \(\Inv_{102}\): \(w_\D=-1/3\) curvature-mimicking degeneracy, exact at
background level.

\item \(\Inv_{103}\): physical \(w_\D=-1/3\) fixed line has \(x_\D>1\) and
\(x_K<0\) in \(S^3\), exact sign correction.

\item \(\Inv_{104}\): \(a_t=(\Omega_{\D0}/|\Omega_K|)^{1/(1+3w_\D)}\) for
\(w_\D\ne-1/3\), exact single-component turning scale.

\item \(\Inv_{105}\): \(a_{\rm acc}=(\Omega_m/2\Omega_\Lambda)^{1/3}\), exact
dust-plus-\(\Lambda\) acceleration threshold.

\item \(\Inv_{106}\): regular FLRW bounce condition
\(H=0\) and \(\epsilon+3p<0\), necessary within classical FLRW GR.

\item \(\Inv_{107}\): \(D_M\le A_0=c/(H_0\sqrt{|\Omega_K|})\), exact closed-geometry
distance bound.

\item \(\Inv_{108}\): interaction scaling ratio
\[
r_*=\frac{-(w_\vis-w_\D)+\sqrt{(w_\vis-w_\D)^2+4\Gamma^2}}{2\Gamma},
\]
conditional on the specified interaction current.

\item \(\Inv_{109}\): Brown--York energy Eq.~\eqref{eq:BY}, exact for the chosen
round slice and reference subtraction.

\item \(\Inv_{110}\): \(E_\BY(A)=c^4A/G\), exact equatorial Brown--York benchmark.

\item \(\Inv_{111}\): Brown--York and Misner--Sharp leading small-sphere agreement
at \(\epsilon=3c^4/(8\pi GA^2)\), exact only through order \(R_A^3\).

\item \(\Inv_{112}\): radius-dependent equality density Eq.~\eqref{eq:eps_BYMS},
exact caution against all-order BY--MS identification.

\item \(\Inv_{113}\): entropy area bound Eq.~\eqref{eq:entropy}, exact safe form.

\item \(\Inv_{114}\): untwisted Dirac spectral gap on round \(S^3\), exact; twisted
gauge-coupled zero-mode absence is not topological.

\item \(\Inv_{115}\): Gauss--Bonnet identity Eq.~\eqref{eq:GB}, exact with the
traceless-Ricci and boundary terms retained.

\item \(\Inv_{116}\): closed-geodesic periods \(L_k=2\pi kA\), exact for round
\(S^3\).

\item \(\Inv_{117}\): Ray--Singer torsion \(T(S^3)=1\) for the trivial real
representation, exact topological benchmark.

\end{enumerate}

%=======================================================================
\section{Summary}
\label{sec:summary}
%=======================================================================

The useful analytic content is not the assertion of a complete solution space.
The useful content is a set of exact benchmark identities against which a compact
\(S^3\) residual-cosmology pipeline can be tested.  The strict core is:

\begin{equation}
\begin{gathered}
\text{proper-time autonomous flow};\\
\text{scalar heat, zeta, and determinant benchmarks};\\
\text{admissible covariance kernels};\\
\text{curvature-mimicking residual backgrounds};\\
\text{acceleration, bounce, and distance criteria};\\
\text{quasi-local energy distinctions}.
\end{gathered}
\end{equation}

Four corrections are especially important for PRD-level rigor.  First, the
physical \(S^3\) curvature variable is negative, so the curvature-mimicking fixed
line has \(x_\D>1\), not \(0\le x_\D\le1\).  Second, the constant-\(w\) turning
scale contains \(1+3w\), not \(1-3w\).  Third, Brown--York and Misner--Sharp
energies agree only at leading small-sphere order under the curvature-density
condition; they are not identical to all orders.  Fourth, index and
Gauss--Bonnet arguments must retain their hypotheses: odd-dimensional Dirac
index arguments do not rule out all twisted zero modes, and four-dimensional
Gauss--Bonnet does not constrain \(\int R^2\) alone for a general FLRW history.

\begin{acknowledgments}
This supplement is designed as a validation document for the benchmark framework.
No numerical or observational claim is made without a corresponding transfer map
and likelihood in the companion pipeline.
\end{acknowledgments}

\begin{thebibliography}{99}

\bibitem{WainwrightEllis1997}
J.~Wainwright and G.~F.~R. Ellis, eds.,
\textit{Dynamical Systems in Cosmology}
(Cambridge University Press, Cambridge, 1997).

\bibitem{Berger1971}
M.~Berger, P.~Gauduchon, and E.~Mazet,
\textit{Le spectre d'une vari\'et\'e riemannienne}
(Springer, Berlin, 1971).

\bibitem{Schoenberg1942}
I.~J. Schoenberg, Positive definite functions on spheres,
Duke Math. J. \textbf{9}, 96 (1942).

\bibitem{Rosenberg1997}
S.~Rosenberg, \textit{The Laplacian on a Riemannian Manifold}
(Cambridge University Press, Cambridge, 1997).

\bibitem{BrownYork1993}
J.~D. Brown and J.~W. York, Jr.,
Quasilocal energy and conserved charges derived from the gravitational action,
Phys. Rev. D \textbf{47}, 1407 (1993).

\bibitem{MisnerSharp1964}
C.~W. Misner and D.~H. Sharp,
Relativistic equations for adiabatic, spherically symmetric gravitational collapse,
Phys. Rev. \textbf{136}, B571 (1964).

\bibitem{Bousso2002}
R.~Bousso, The holographic principle,
Rev. Mod. Phys. \textbf{74}, 825 (2002).

\bibitem{Wald1984}
R.~M. Wald, \textit{General Relativity}
(University of Chicago Press, Chicago, 1984).

\bibitem{Lichnerowicz1958}
A.~Lichnerowicz,
\textit{G\'eom\'etrie des groupes de transformations}
(Dunod, Paris, 1958).

\bibitem{RaySinger1971}
D.~B. Ray and I.~M. Singer,
\(R\)-torsion and the Laplacian on Riemannian manifolds,
Adv. Math. \textbf{7}, 145 (1971).

\end{thebibliography}

\end{document}
